% !TEX root = ../algebra_02.tex

\section{Факторгруппы, гомоморфизмы и циклические группы}

\subsection{Факторгруппы}

\begin{Definition}{Фактормножество}{def:quotient-set}
	Пусть на множестве $X$ задано отношение эквивалентности $\sim$. Тогда \emph{фактормножество} $X/{\sim}$ определяется как множество всех классов эквивалентности:
	\[
	X/{\sim}:=\{[x] \mid x\in X\}.
	\]
\end{Definition}

\begin{Remark}{Фактормножество группы по нормальной подгруппе}{rem:group-quotient-set}
	Пусть $G$ --- группа, $H \triangleleft G$. Тогда на $G$ задано отношение эквивалентности
	\[
	g_1\sim g_2 \Longleftrightarrow g_1^{-1}g_2\in H.
	\]
	Его классы эквивалентности --- это левые смежные классы по $H$, поэтому соответствующее фактормножество имеет вид
	\[
	G/H:=\{gH\mid g\in G\}.
	\]
	На $G/H$ задаётся операция
	\[
	(gH)\cdot(g'H):=(gg')H.
	\]
	Так как $H$ нормальна, левые и правые смежные классы совпадают: $gH=Hg$ для всех $g\in G$.
\end{Remark}

\begin{Proposition}{Корректность умножения смежных классов}{prop:well-defined}
	Если $H \triangleleft G$, то операция
	\[
	(gH)(g'H) := (gg')H
	\]
	корректно определена, то есть не зависит от выбора представителей смежных классов.
\end{Proposition}

\begin{proof}
	Пусть $g_1H=g_2H, \qquad g_1'H=g_2'H$.
	
	Тогда существуют $h,h'\in H$ такие, что $g_2=g_1h, \qquad g_2'=g_1'h'$.

	Следовательно, $g_2g_2' = g_1 h g_1' h' = g_1 g_1'\bigl((g_1')^{-1} h g_1'\bigr)h'$.

	Так как $H \triangleleft G$, имеем $(g_1')^{-1} h g_1' \in H$. Поэтому $\bigl((g_1')^{-1} h g_1'\bigr)h' \in H,$ и значит, $g_2g_2'H = g_1g_1'H.$

	Следовательно, $(g_2H)(g_2'H)=(g_1H)(g_1'H).$

	Корректность доказана.
\end{proof}

\begin{Theorem}{Групповая структура на $G/H$}{thm:quotient-group}
	Если $H \triangleleft G$, то множество $G/H$ с операцией
	\[
	(gH)(g'H)=(gg')H
	\]
	является группой.
\end{Theorem}

\begin{proof}
	Проверим аксиомы группы.
	\begin{enumerate}
		\item Ассоциативность:
		\[
		\bigl((g_1H)(g_2H)\bigr)(g_3H) = (g_1g_2g_3)H = (g_1H)\bigl((g_2H)(g_3H)\bigr).
		\]
		\item Нейтральный элемент: $eH=H$, и для любого $g\in G$
		\[
		(eH)(gH)=gH=(gH)(eH).
		\]
		\item Обратимый элемент: для $gH\in G/H$
		\[
		(gH)(g^{-1}H)=eH=(g^{-1}H)(gH).
		\]
	\end{enumerate}
	Следовательно, $G/H$ --- группа.
\end{proof}

По определению эта группа называется \emph{факторгруппой} группы $G$ по нормальной подгруппе $H$.

\begin{Remark}{}{rem:not-subgroup}
	Вообще говоря, $G/H$ не является подмножеством группы $G$. Это новая группа, элементы которой --- смежные классы. Исключения возникают лишь в тривиальных случаях $H=\{e\}$ и $H=G$.
\end{Remark}

\subsection{Гомоморфизмы групп}

\begin{Definition}{Гомоморфизм групп}{def:hom}
	Пусть $G$ и $G'$ --- группы. Отображение $\varphi\colon G \to G'$
	называется \emph{гомоморфизмом групп}, если для всех $g_1,g_2\in G$ выполнено
	\[
	\varphi(g_1g_2)=\varphi(g_1)\varphi(g_2).
	\]
\end{Definition}

\begin{Example}{}{ex:hom-examples}
	\begin{enumerate}
		\item Определитель:
		\[
		\det\colon \GL_n(K) \to K^{\times}, \qquad \det(AB)=\det(A)\det(B).
		\]
		\item Знак перестановки:
		\[
		\sgn\colon S_n \to \{\pm 1\}, \qquad \sigma \mapsto \sgn(\sigma).
		\]
		\item Тождественный и нулевой гомоморфизмы:
		\[
		\id_G\colon G\to G, \quad g\mapsto g;
		\qquad
		0\colon G\to G', \quad g\mapsto e_{G'}.
		\]
		\item Логарифм как гомоморфизм между мультипликативной и аддитивной группами:
		\[
		\ln\colon (\RR_{>0},\cdot) \to (\RR,+), \qquad a\mapsto \ln a.
		\]
	\end{enumerate}
\end{Example}

\begin{Lemma}{Сохранение единицы и обратного элемента}{lem:unit-inverse}
	Пусть $\varphi\colon G\to G'$ --- гомоморфизм. Тогда
	\begin{enumerate}
		\item $\varphi(e_G)=e_{G'}$;
		\item для любого $g\in G$ выполнено
		\[
		\varphi(g^{-1})=\varphi(g)^{-1}.
		\]
	\end{enumerate}
\end{Lemma}

\begin{proof}
	Так как $e_Ge_G=e_G$, имеем
	\[
	\varphi(e_G)\varphi(e_G)=\varphi(e_G).
	\]
	Умножая это равенство слева на $\varphi(e_G)^{-1}$, получаем $\varphi(e_G)=e_{G'}$.
	
	Далее,
	\[
	\varphi(g^{-1})\varphi(g)=\varphi(g^{-1}g)=\varphi(e_G)=e_{G'},
	\]
	и аналогично
	\[
	\varphi(g)\varphi(g^{-1})=e_{G'}.
	\]
	Следовательно, $\varphi(g^{-1})=\varphi(g)^{-1}$.
\end{proof}

\begin{Proposition}{Образы и прообразы подгрупп}{image-preimage}
	Пусть $\varphi\colon G\to G'$ --- гомоморфизм.
	\begin{enumerate}
		\item Если $H\le G$, то $\varphi(H)\le G'$.
		\item Если $H'\le G'$, то $\varphi^{-1}(H')\le G$.
		\item Если $H'\triangleleft G'$, то $\varphi^{-1}(H')\triangleleft G$.
	\end{enumerate}
\end{Proposition}

\begin{proof}
	\begin{enumerate}
		\item Пусть $g_1',g_2'\in \varphi(H)$. Тогда существуют $g_1,g_2\in H$ такие, что $g_1'=\varphi(g_1)$ и $g_2'=\varphi(g_2)$. Поскольку $H$ --- подгруппа, $g_1g_2\in H, \qquad g_1^{-1}\in H.$\\
		Следовательно, $g_1'g_2' = \varphi(g_1)\varphi(g_2)=\varphi(g_1g_2)\in \varphi(H),$ и по лемме $(g_1')^{-1}=\varphi(g_1)^{-1}=\varphi(g_1^{-1})\in \varphi(H).$\\
		Значит, $\varphi(H)\le G'$.
		
		\item Пусть $g_1,g_2\in \varphi^{-1}(H')$. Тогда $\varphi(g_1),\varphi(g_2)\in H'$.\\
		Так как $H'$ --- подгруппа, $\varphi(g_1g_2)=\varphi(g_1)\varphi(g_2)\in H',$ то есть $g_1g_2\in \varphi^{-1}(H')$.\\
		Кроме того, $\varphi(g_1^{-1})=\varphi(g_1)^{-1}\in H',$ поэтому $g_1^{-1}\in \varphi^{-1}(H')$. Значит, $\varphi^{-1}(H')\le G$.
		
		\item Пусть $g\in G$ и $h\in \varphi^{-1}(H')$. Тогда $\varphi(h)\in H'$. Поскольку $H'\triangleleft G'$, имеем
		\[
		\varphi(ghg^{-1})=\varphi(g)\varphi(h)\varphi(g)^{-1}\in H'.
		\]
		Следовательно, $ghg^{-1}\in \varphi^{-1}(H')$. Значит, $\varphi^{-1}(H')\triangleleft G$.
	\end{enumerate}
\end{proof}

\begin{Corollary}{Образ и ядро}{prop:image-kernel}
	Для любого гомоморфизма $\varphi\colon G\to G'$ имеем
	\[
	\Im \varphi := \varphi(G) \le G',
	\qquad
	\Ker \varphi := \varphi^{-1}(\{e_{G'}\}) = \{g\in G \mid \varphi(g)=e_{G'}\} \triangleleft G.
	\]
\end{Corollary}

\begin{proof}
	Первая часть следует из пункта 1) предложения~\ref{prop:image-preimage}, вторая --- из пункта 3) того же предложения, поскольку $\{e_{G'}\} \triangleleft G'$.
\end{proof}

\begin{Example}{}{ex:alternating}
	Имеем $A_n = \Ker(\sgn) \triangleleft S_n.$
	
	Поэтому знак перестановки даёт изоморфизм
	\[
	S_n/A_n \cong \{\pm 1\}.
	\]
\end{Example}

\begin{Proposition}{Критерий инъективности}{prop:inj-kernel}
	Пусть $\varphi\colon G\to G'$ --- гомоморфизм. Тогда
	\[
	\varphi \text{ инъективен } \Longleftrightarrow \Ker\varphi = \{e_G\}.
	\]
\end{Proposition}

\begin{proof}
	Предположим, что $\varphi$ инъективен. 
	
	Если $h\in \Ker\varphi$, то $\varphi(h)=e_{G'}=\varphi(e_G).$
	
	Из инъективности следует $h=e_G$, значит, $\Ker\varphi=\{e_G\}$.
	
	Обратно, пусть $\Ker\varphi=\{e_G\}$ и $\varphi(h_1)=\varphi(h_2).$

	Тогда $\varphi(h_1h_2^{-1})=\varphi(h_1)\varphi(h_2)^{-1}=e_{G'},$ то есть $h_1h_2^{-1}\in \Ker\varphi$. Следовательно, $h_1h_2^{-1}=e_G,$ откуда $h_1=h_2$. Значит, $\varphi$ инъективен.
\end{proof}

\begin{Definition}{Специальные виды гомоморфизмов}{def:special-homs}
	\begin{enumerate}
		\item Инъективный гомоморфизм называется \emph{мономорфизмом}.
		\item Сюръективный гомоморфизм называется \emph{эпиморфизмом}.
		\item Биективный гомоморфизм называется \emph{изоморфизмом}.
		\item Гомоморфизм $G\to G$ называется \emph{эндоморфизмом}.
		\item Изоморфизм $G\to G$ называется \emph{автоморфизмом}.
	\end{enumerate}
\end{Definition}

\begin{Proposition}{Обратное отображение к изоморфизму}{prop:inverse-iso}
	Если $\varphi\colon G\to G'$ --- изоморфизм, то обратное отображение
	\[
	\varphi^{-1}\colon G'\to G
	\]
	тоже является изоморфизмом.
\end{Proposition}

\begin{proof}
	Так как $\varphi$ биективно, отображение $\varphi^{-1}$ существует и биективно. Для любых $g_1',g_2'\in G'$ положим
	\[
	A:=\varphi^{-1}(g_1'g_2'), \qquad B:=\varphi^{-1}(g_1')\varphi^{-1}(g_2').
	\]
	Тогда
	\[
	\varphi(A)=g_1'g_2'=
	\varphi\bigl(\varphi^{-1}(g_1')\bigr)\,\varphi\bigl(\varphi^{-1}(g_2')\bigr)=\varphi(B).
	\]
	Из инъективности $\varphi$ следует $A=B$. Следовательно,
	\[
	\varphi^{-1}(g_1'g_2')=\varphi^{-1}(g_1')\varphi^{-1}(g_2').
	\]
	Значит, $\varphi^{-1}$ --- гомоморфизм, а потому изоморфизм.
\end{proof}

\begin{Example}{Простейшие примеры факторгрупп}{ex:trivial-quotients}
	Имеют место изоморфизмы
	\[
	G/\{e\} \cong G, \qquad G/G \cong \{e\}.
	\]
	Если группа записана аддитивно, то факторгруппа тоже записывается аддитивно. Например,
	\[
	\ZZ/m\ZZ = \{[a]_m \mid a\in \ZZ\}, \qquad [a]_m + [b]_m := [a+b]_m.
	\]
\end{Example}

\subsection{Канонические гомоморфизмы и индуцированные отображения}

\begin{Definition}{Канонические гомоморфизмы}{def:canonical-homs}
	Пусть $H\le G$.
	\begin{enumerate}
		\item \textbf{Гомоморфизм включения:}
		\[
		i_H\colon H\to G, \qquad h\mapsto h.
		\]
		Он инъективен, причём
		\[
		\Im i_H = H, \qquad \Ker i_H = \{e\}.
		\]
		\item Если дополнительно $H\triangleleft G$, то \textbf{канонический эпиморфизм} (проекция) задаётся формулой
		\[
		\pi_H\colon G\to G/H, \qquad g\mapsto gH.
		\]
		Это сюръективный гомоморфизм, причём
		\[
		\Im \pi_H = G/H, \qquad \Ker \pi_H = H.
		\]
	\end{enumerate}
\end{Definition}

\begin{proof}
	Гомоморфность проекции следует из равенства
	\[
	\pi_H(g_1g_2)=g_1g_2H=(g_1H)(g_2H)=\pi_H(g_1)\pi_H(g_2).
	\]
	Кроме того,
	\[
	\pi_H(g)=eH \Longleftrightarrow gH=H \Longleftrightarrow g\in H.
	\]
	Значит, $\Ker \pi_H=H$.
\end{proof}

\begin{Theorem}{Индуцированный гомоморфизм}{induced-hom}
	Пусть дан гомоморфизм
	\[
	\varphi\colon G\to G'.
	\]
	Пусть $H\triangleleft G$ и $H\subseteq \Ker\varphi$. Пусть также $H'\le G'$ и
	$\Im\varphi\subseteq H'$. Тогда отображение
	\[
	\widetilde\varphi\colon G/H\to H', \qquad \widetilde\varphi(gH):=\varphi(g),
	\]
	корректно определено и является гомоморфизмом. Более того, это единственный гомоморфизм для которого коммутирует диаграмма
	\[
	\begin{tikzcd}[column sep=large, row sep=large]
		G \arrow[r, "\varphi"] \arrow[d, "\pi_H"'] & G' \\
		G/H \arrow[r, "\widetilde{\varphi}"'] & H' \arrow[u, "i_{H'}"']
	\end{tikzcd}
	\]
	то есть
	\[
	\varphi = i_{H'}\circ \widetilde\varphi\circ \pi_H.
	\]
\end{Theorem}

\begin{proof}
	Сначала докажем корректность.
	Если $g_1H=g_2H$, то $g_2=g_1h$ для некоторого $h\in H$. Так как
	$H\subseteq \Ker\varphi$, имеем $\varphi(h)=e$, и потому
	\[
	\varphi(g_2)=\varphi(g_1h)=\varphi(g_1)\varphi(h)=\varphi(g_1).
	\]
	Значит, значение $\widetilde\varphi(gH)$ не зависит от выбора представителя класса.
	
	Так как $\Im\varphi\subseteq H'$, действительно $\widetilde\varphi(gH)=\varphi(g)\in H',$ то есть отображение $\widetilde\varphi$ принимает значения в $H'$.
	
	Проверим, что $\widetilde\varphi$ является гомоморфизмом:
	\[
	\widetilde\varphi\bigl((g_1H)(g_2H)\bigr)
	=
	\widetilde\varphi(g_1g_2H)
	=
	\varphi(g_1g_2)
	=
	\varphi(g_1)\varphi(g_2)
	=
	\widetilde\varphi(g_1H)\widetilde\varphi(g_2H).
	\]
	
	Теперь проверим коммутативность диаграммы. Для любого $g\in G$ имеем
	\[
	(i_{H'}\circ \widetilde\varphi\circ \pi_H)(g)
	=
	i_{H'}\bigl(\widetilde\varphi(gH)\bigr)
	=
	i_{H'}(\varphi(g))
	=
	\varphi(g).
	\]
	Следовательно, 	$\varphi = i_{H'}\circ \widetilde\varphi\circ \pi_H.$
	
	Докажем единственность. Пусть $\psi\colon G/H\to H'$ --- гомоморфизм такой, что $\varphi = i_{H'}\circ \psi\circ \pi_H.$
	
	Тогда для любого $g\in G$
	\[
	i_{H'}\bigl(\psi(gH)\bigr)
	=
	(i_{H'}\circ \psi\circ \pi_H)(g)
	=
	\varphi(g)
	=
	i_{H'}\bigl(\widetilde\varphi(gH)\bigr).
	\]
	Так как $i_{H'}\colon H'\hookrightarrow G'$ --- вложение, оно инъективно. Значит,
	\[
	\psi(gH)=\widetilde\varphi(gH)
	\]
	для всех $g\in G$, откуда $\psi=\widetilde\varphi$.
\end{proof}

\begin{Theorem}{Первая теорема об изоморфизме}{thm:first-iso}
	Пусть $\varphi\colon G\to G'$ --- гомоморфизм. Тогда индуцированный гомоморфизм
	\[
	\overline\varphi\colon G/\Ker\varphi \to \Im\varphi,
	\qquad
	\overline\varphi(g\Ker\varphi)=\varphi(g)
	\]
	является изоморфизмом. В частности,
	\[
	G/\Ker\varphi \cong \Im\varphi.
	\]
\end{Theorem}

\begin{proof}
	Существование и гомоморфность следуют из теоремы~\ref{thm:induced-hom}. По определению
	\[
	\Im\overline\varphi=\Im\varphi.
	\]
	Кроме того,
	\[
	\Ker\overline\varphi = \{g\Ker\varphi \mid \varphi(g)=e\} = \{g\Ker\varphi \mid g\in \Ker\varphi\} = \{\Ker\varphi\}.
	\]
	Следовательно, $\overline\varphi$ инъективен. Значит, $\overline\varphi$ --- изоморфизм.
\end{proof}

\begin{Example}{Применение первой теоремы об изоморфизме}{ex:sln}
	Определитель задаёт сюръективный гомоморфизм
	\[
	\det\colon \GL_n(K)\to K^{\times}.
	\]
	Его ядро равно
	\[
	\Ker(\det)=\SL_n(K)=\{A\in \GL_n(K)\mid \det A=1\}.
	\]
	Поэтому
	\[
	\SL_n(K)\triangleleft \GL_n(K),
	\qquad
	\GL_n(K)/\SL_n(K)\cong K^{\times}.
	\]
\end{Example}

\subsection{Циклические группы}

\begin{Definition}{Циклическая группа}{def:cyclic}
	Группа $G$ называется \emph{циклической}, если существует элемент $g\in G$ такой, что
	\[
	G=\langle g\rangle = \{g^k\mid k\in \ZZ\}.
	\]
	Элемент $g$ называется \emph{образующим} группы $G$.
\end{Definition}

\begin{Theorem}{Все подгруппы группы $\ZZ$}{subgroups-Z}
	Всякая подгруппа $H\le \ZZ$ имеет вид
	\[
	H=n\ZZ=\langle n\rangle
	\]
	для некоторого $n\ge 0$.
\end{Theorem}

\begin{proof}
	Если $H=\{0\}$, то утверждение верно при $n=0$.
	
	Пусть теперь $H\ne\{0\}$. Тогда в $H$ есть положительные элементы. Пусть $n$ --- наименьший положительный элемент из $H$. Покажем, что $H=n\ZZ$.
	
	Во-первых, $n\ZZ\subseteq H$, поскольку $H$ замкнута относительно сложения и взятия противоположного.
	
	Во-вторых, пусть $x\in H$. По алгоритму деления существует представление
	\[
	x=qn+r, \qquad 0\le r<n.
	\]
	Так как $x\in H$ и $qn\in H$, получаем $r=x-qn\in H.$
	По минимальности $n$ имеем $r=0$. Значит, $x\in n\ZZ$.
	
	Следовательно, $H=n\ZZ$.
\end{proof}

\begin{Theorem}{Классификация циклических групп}{thm:cyclic-classification}
	Пусть $G=\langle g\rangle$ --- циклическая группа.
	\begin{enumerate}
		\item Если $|G|=\infty$, то $G\cong \ZZ$.
		\item Если $|G|=n<\infty$, то $G\cong \ZZ/n\ZZ$.
	\end{enumerate}
\end{Theorem}

\begin{proof}
	Рассмотрим гомоморфизм
	\[
	\varphi\colon \ZZ \to G, \qquad k\mapsto g^k.
	\]
	Он сюръективен, поскольку $G=\langle g\rangle$. По теореме~\ref{thm:subgroups-Z} существует $n\ge 0$ такое, что $\Ker\varphi=n\ZZ.$

	Если $n=0$, то $\Ker\varphi=\{0\}$, и потому $\varphi$ инъективен. Следовательно, $G\cong \ZZ$.
	
	Если $n>0$, то по первой теореме об изоморфизме
	\[
	\ZZ/n\ZZ = \ZZ/\Ker\varphi \cong \Im\varphi = G.
	\]
	Если при этом $|G|<\infty$, то обязательно $n=|G|$.
\end{proof}

\subsection{Теоремы соответствия и об изоморфизме}

\begin{Theorem}{Теорема соответствия}{thm:correspondence}
	Пусть $H\triangleleft G$ и $\pi_H\colon G\to G/H$ --- каноническая проекция. Тогда отображения
	\[
	\alpha\colon \{P\le G\mid H\subseteq P\}\to \{Q\le G/H\}, \qquad \alpha(P)=\pi_H(P),
	\]
	\[
	\beta\colon \{Q\le G/H\}\to \{P\le G\mid H\subseteq P\}, \qquad \beta(Q)=\pi_H^{-1}(Q)
	\]
	взаимно обратны. Следовательно, между подгруппами группы $G$, содержащими $H$, и подгруппами факторгруппы $G/H$ существует биекция.
\end{Theorem}

\begin{proof}
	Покажем, что $\alpha$ и $\beta$ корректны и взаимно обратны.
	
	Если $P\le G$ и $H\subseteq P$, то $\pi_H(P)$ является подгруппой группы $G/H$, поскольку это образ подгруппы $P$ при гомоморфизме $\pi_H|_P\colon P\to G/H$.
	
	Если $Q\le G/H$, то $\pi_H^{-1}(Q)\le G$ по предложению~\ref{prop:image-preimage}; кроме того, $H=\pi_H^{-1}(\{eH\})\subseteq \pi_H^{-1}(Q)$.
	
	Теперь вычислим композиции.
	
	Для $Q\le G/H$ имеем
	\[
	(\alpha\circ\beta)(Q)=\alpha\bigl(\pi_H^{-1}(Q)\bigr)=\pi_H\bigl(\pi_H^{-1}(Q)\bigr)=Q,
	\]
	поскольку $\pi_H$ сюръективна.
	
	Для $P\le G$ при $H\subseteq P$ имеем
	\[
	(\beta\circ\alpha)(P)=\pi_H^{-1}(\pi_H(P))=\{g\in G\mid gH\in \pi_H(P)\}.
	\]
	Но условие $gH\in \pi_H(P)$ означает, что существует $p\in P$ такое, что $gH=pH$. Тогда $p^{-1}g\in H\subseteq P$, значит, $g\in P$. Итак,
	\[
	\pi_H^{-1}(\pi_H(P))=P.
	\]
	Следовательно, $\alpha$ и $\beta$ взаимно обратны.
\end{proof}

\begin{Statement}{нормальность в факторгруппе}{normality-correspondence}
	Пусть $H\triangleleft G$, $P\le G$ и $H\subseteq P$. Тогда
	\[
	\pi_H(P) \triangleleft G/H \Longleftrightarrow P \triangleleft G.
	\]
\end{Statement}

\begin{proof}
	Поскольку $\pi_H\colon G\to G/H$ --- сюръективный гомоморфизм и
	\[
	\pi_H^{-1}(\pi_H(P))=P,
	\]
	одна импликация сразу следует из свойства прообраза нормальной подгруппы:
	если $\pi_H(P)\triangleleft G/H$, то
	\[
	P=\pi_H^{-1}(\pi_H(P))\triangleleft G.
	\]
	
	Обратно, пусть $P\triangleleft G$. Рассмотрим ограничение проекции на $P$. Тогда $\pi_H(P)$ --- подгруппа в $G/H$ по теореме соответствия. Для любого $gH\in G/H$ и любого $pH\in \pi_H(P)$ имеем
	\[
	(gH)(pH)(gH)^{-1}=(gpg^{-1})H.
	\]
	Так как $P\triangleleft G$, элемент $gpg^{-1}$ принадлежит $P$. Значит,
	\[
	(gpg^{-1})H\in \pi_H(P).
	\]
	Следовательно, $\pi_H(P)\triangleleft G/H$.
\end{proof}

\begin{Theorem}{Третья теорема об изоморфизме}{thm:third-iso}
	Пусть $H\triangleleft G$, $P\triangleleft G$ и $H\subseteq P$. Тогда
	\begin{enumerate}
		\item $H\triangleleft P$;
		\item $P/H\triangleleft G/H$;
		\item $(G/H)/(P/H) \cong G/P.$

	\end{enumerate}
\end{Theorem}